\documentclass[
12pt,
BCOR=10mm,
a4paper,
bibliography=totoc,
includefoot,
listof=totoc,
smallheadings,
]{scrartcl}

% charset, encoding and default language
\usepackage[T1]{fontenc}
\usepackage[utf8]{inputenc}
\usepackage[ngerman]{babel}

% use helvetica
% set sans serif as default font family
\usepackage[scaled]{helvet}
\renewcommand{\familydefault}{\sfdefault}

% improve general appearance of type
\usepackage{microtype}

% line height
\usepackage{setspace}
\onehalfspacing

% recalculate type area
\typearea[current]{calc}

% custom page layout
\usepackage[
a4paper,
bottom=3cm,
left=3cm,
right=3cm,
top=2.5cm,
]{geometry}

% custom header and footer
\usepackage[
autooneside=false,
headsepline,
]{scrlayer-scrpage}

% default page style
\pagestyle{scrheadings}

% save roman page number for later use
\newcounter{savepage}

% use these titles
\automark[subsection]{section}

% clean and set marks
\chead{}
% \ohead[\headmark]{\headmark}
\cfoot{}
\ofoot[\pagemark]{\pagemark}

% colorize
\usepackage{xcolor}
\setkomafont{pageheadfoot}{%
    \itshape\color{gray}}
\setkomafont{pagenumber}{%
    \itshape\color{gray}}

% graphics
\usepackage{graphicx}
% tables
\usepackage{booktabs}

% papers
\usepackage[backend=biber]{biblatex}
\addbibresource{papers.bib}

% abbreviations
\usepackage[intoc]{nomencl}
\renewcommand{\nomname}{Abkürzungsverzeichnis}
\renewcommand{\nomlabel}[1]{\textbf{#1}}
\makenomenclature

% start sections on a new page
\usepackage{titlesec}
\newcommand{\sectionbreak}{\clearpage}

% adjust font size of headings
\titleformat*{\section}{\Large\bfseries}
\titleformat*{\subsection}{\large\bfseries}
\titleformat*{\subsubsection}{\normalsize\bfseries}
% \titleformat*{\paragraph}{\normalsize\bfseries}
% \titleformat*{\subparagraph}{\normalsize\bfseries}

% nice references and links
\usepackage[german]{fancyref}
\usepackage[pdfusetitle]{hyperref}

% lorem ipsum
\usepackage{lipsum}

% here we go
\begin{document}

\title{My Awesome Thesis}
\author{Philipp Kosel}

\maketitle
\newpage

% breakpoint 1
\cleardoublepage
\pagenumbering{Roman}
\ohead[]{}

\addsec{Kurzfassung}
\lipsum[1]
\newpage

\addsec{Abstract}
\lipsum[1]
\newpage

\tableofcontents
\newpage

\listoffigures
\newpage

\listoftables
\newpage

\printnomenclature
\newpage

% breakpoint 2
\cleardoublepage
\setcounter{savepage}{\arabic{page}}
\pagenumbering{arabic}
\ohead[\headmark]{\headmark}

% ch01
\section{Headline Level 1}
\lipsum[1-2]

Lorem ipsum dolor sid amet \autocite{mcconnell2004}.

Die Abkürzung DHBW\nomenclature{DHBW}{Duale Hochschule Baden Württemberg} steht im Abkürzungsverzeichnis.

Verweis auf~\fref{fig:logo}.

\subsection{Headline Level 2}
\lipsum[1]

\begin{table}[htbp]
    \caption{Eine Tabelle}
    \label{tab:tabelle}
    \centering
    \begin{tabular}{lll}
        \toprule
        Ziel      & Abfahrt   & Dauer \\
        \midrule
        Frankfurt & stündlich & 3:20  \\
        Berlin    & stündlich & 5:40  \\
        Hamburg   & alle 5\,h & 5:20  \\
        \bottomrule
    \end{tabular}
\end{table}

\subsection{Headline Level 2}
\lipsum[1]

\begin{figure}[htbp]
    \centering
    \includegraphics[scale=1.6]{figures/DHBW_d_LOE_46mm_4c.pdf}
    \caption{Das Logo der DHBW Lörrach}
    \label{fig:logo}
\end{figure}

% ch02
\section{Headline Level 1}
\lipsum[1]
\subsection{Headline Level 2}
\lipsum[1]
\subsubsection{Headline Level 3}
\lipsum[1]
\subsubsection{Headline Level 3}
\lipsum[1]
\subsection{Headline Level 2}
\lipsum[1]

% ch03
\section{Headline Level 1}
\lipsum[1]
\subsection{Headline Level 2}
\lipsum[1]
\subsection{Headline Level 2}
\lipsum[1]

% ch04
\section{Headline Level 1}
\lipsum[1]
\subsection{Headline Level 2}
\lipsum[1]
\subsection{Headline Level 2}
\lipsum[1]

% ch05
\section{Headline Level 1}
\lipsum[1]
\subsection{Headline Level 2}
\lipsum[1]
\subsection{Headline Level 2}
\lipsum[1]

\newpage

% breakpoint 3
\cleardoublepage
\pagenumbering{Roman}
\setcounter{page}{\thesavepage}
\ohead[]{}

\printbibliography

\end{document}
